% Draft subsection: an alternative argument for the orthogonal coercivity bound
% via weighted Bochner identities, and an improved Cauchy--Schwarz inequality
% for entire functions vanishing at the origin.
%
% Notation matches spr-hermite.tex (Fock measure
% \dd\gamma(z) = (1/\pi)\ee^{-\abs{z}^2}\dd m(z), Fock norm \norm{\cdot}_{\mathcal{F}},
% \abs{\cdot}, \norm{\cdot}, \rho(w) = \abs{\abs{1+w}-1}, \dd, \ee, \ii, \C)
% and fnt_logsob_proof.tex (\operatorname{Re}, \operatorname{Im},
% \langle\cdot,\cdot\rangle_{\mathcal{F}}).
%
% Suggested placement: as a new subsection in Section~\ref{subsec:assembly}'s
% parent, e.g.\ between \subsection{Proof of the orthogonal coercivity estimate}
% (Section~\ref{subsec:assembly}) and \subsection{Generalizations of
% Theorem~\ref{main:local}}.

\subsection{An alternative approach via weighted Bochner identities}\label{subsec:bochner}

In this subsection we present a complementary route to the orthogonal
coercivity bound of Theorem~\ref{thm:orthog} (and hence to
Theorem~\ref{main:local}). The argument bypasses the circle estimates
of Section~\ref{subsec:circle} and the frequency localization of
Section~\ref{subsec:blocks}, replacing them with a direct integration
by parts in the polynomial weight $(1+\abs{z}^2)^{-1}\,\dd\gamma$ together
with a weighted Cauchy integral inequality. As a byproduct we obtain a
quantitative \emph{improved Cauchy--Schwarz inequality} between
$\operatorname{Re}(G)$ and $\abs{G}^2$ in $L^2(\dd\gamma)$ for entire
functions $G$ vanishing at the origin, in the spirit of Aldaz's
sharpening of Cauchy--Schwarz.

\medskip

Throughout this subsection $F, G \in \mathcal{F}^2(\C)$ denote polynomials
unless stated otherwise; the conclusions then extend to arbitrary
elements of $\mathcal{F}^2(\C)$ by a standard density argument analogous
to the one in Subsection~\ref{subsec:local-reduction}. We also abbreviate
\[
\mathcal{F}^2_\ast(\C) \;:=\; \bigl\{H \in \mathcal{F}^2(\C) : H(0) = 0\bigr\},
\]
the closed subspace of $\mathcal{F}^2(\C)$ orthogonal to the constant
function~$\mathbf{1}$.

\subsubsection{A weighted equivalent norm}\label{sec:weighted-norm}

The starting point is a weighted Hardy-type identity that expresses the
Fock norm of an element of $\mathcal{F}^2_\ast(\C)$ in terms of a polynomial-weighted
$L^2$ integral of its derivative.

\begin{lemma}\label{lem:weighted-equiv}
There exist absolute constants $c_1, c_2 \in (0,\infty)$ such that for
every $H \in \mathcal{F}^2_\ast(\C)$,
\begin{equation}\label{eq:weighted-equiv}
c_1\, \norm{H}_{\mathcal{F}}^2
\;\le\;
\int_{\C} \abs{H'(z)}^2\, \frac{\dd\gamma(z)}{1+\abs{z}^2}
\;\le\;
c_2\, \norm{H}_{\mathcal{F}}^2.
\end{equation}
\end{lemma}

\begin{proof}
Expand $H(z) = \sum_{n \ge 1} a_n z^n$ with $\sum_{n\ge1} \abs{a_n}^2 n!
= \norm{H}_{\mathcal{F}}^2$. Since the monomials $z^n$ are pairwise
orthogonal in $L^2(\dd\gamma)$, the integral on the left-hand side of
\eqref{eq:weighted-equiv} decomposes as
\[
\int_{\C} \abs{H'(z)}^2 \frac{\dd\gamma(z)}{1+\abs{z}^2}
= \sum_{n \ge 1} n^2\, \abs{a_n}^2\, \kappa_n,
\quad
\kappa_n := \int_{\C} \frac{\abs{z}^{2(n-1)}}{1+\abs{z}^2}\, \dd\gamma(z).
\]
Writing $1/(1+\abs{z}^2) = \int_0^\infty \ee^{-(1+\abs{z}^2)s}\,\dd s$ and
using the Gaussian moment $\int_{\C}\abs{z}^{2k}\,\dd\gamma = k!$,
\[
\kappa_n
= (n-1)! \int_0^\infty \frac{\ee^{-s}}{(1+s)^n}\, \dd s.
\]
The integral $\int_0^\infty \ee^{-s}(1+s)^{-n}\,\dd s$ is positive, bounded
above by $\int_0^\infty (1+s)^{-n}\dd s = 1/(n-1)$ for $n \ge 2$, and
satisfies $\int_0^\infty \ee^{-s}(1+s)^{-n}\dd s = (1/n) + O(1/n^2)$ for
$n \to \infty$ (by Watson's lemma at $s=0$). Hence
$n^2\,\kappa_n / n! = n \int_0^\infty \ee^{-s}(1+s)^{-n}\dd s \in [c_1, c_2]$
for some absolute constants $0 < c_1 \le c_2 < \infty$, and
\eqref{eq:weighted-equiv} follows by summing in~$n$.
\end{proof}

\begin{remark}
A short calculation gives the explicit values $c_1 = 1 - \ee\, E_1(1)
\approx 0.404$ (attained at $n=1$) and $c_2 \le 2$, so the
two-sided estimate~\eqref{eq:weighted-equiv} can be made effective. In
the sequel, however, only the qualitative equivalence is used.
\end{remark}

\subsubsection{Stability of the squaring map}\label{sec:squaring}

We now state the main estimate of this subsection, which controls the
Fock norm of $F^2 - F(0)^2$ by the $L^2(\dd\gamma)$ distance from
$\abs{F}^2$ to a real harmonic function.

\begin{theorem}\label{thm:squaring}
There exists an absolute constant $C > 0$ such that for every polynomial
$F \in \mathcal{F}^2(\C)$ and every real harmonic function $v$ on $\C$
satisfying $\abs{F}^2 - v \in L^2(\dd\gamma)$,
\[
\norm{F^2 - F(0)^2}_{\mathcal{F}}
\;\le\;
C\, \bigl\|\, \abs{F}^2 - v \,\bigr\|_{L^2(\dd\gamma)}.
\]
In particular, taking $v = \int_{\C}\abs{F}^2\,\dd\gamma$ (a constant, hence
harmonic),
\begin{equation}\label{eq:squaring-min}
\norm{F^2 - F(0)^2}_{\mathcal{F}}
\;\le\;
C\, \min_{c \in \R}\bigl\|\, \abs{F}^2 - c \,\bigr\|_{L^2(\dd\gamma)}.
\end{equation}
\end{theorem}

The proof occupies the rest of this sub-subsection. We follow the strategy
of integrating by parts twice in the weighted measure
$(1+\abs{z}^2)^{-2}\,\dd\gamma$, the second time after replacing $\abs{F}^2$
by $\abs{F}^2 - v$, which is permissible since $v$ is harmonic and only
the Laplacian of $\abs{F}^2$ enters.

\medskip
\noindent\emph{Step 1: Cauchy--Riemann reduction.}\;
Recall from~\eqref{eq:fnt-nabla-mod} that, for an entire function~$F$,
$\abs{\nabla \abs{F}}^2 = \abs{F'}^2$ a.e.\ on $\C$. Multiplying by
$\abs{F}^2$ on both sides one obtains, again from the Cauchy--Riemann
equations,
\begin{equation}\label{eq:bochner-pointwise}
\abs{\nabla \abs{F}^2}^2
\;=\;
4\, \abs{F}^2\, \abs{F'}^2
\;=\;
\abs{(F^2)'}^2
\qquad\text{a.e.\ on } \C.
\end{equation}
Indeed $(F^2)' = 2F F'$, so $\abs{(F^2)'}^2 = 4\abs{F}^2 \abs{F'}^2$, while
writing $F = u + \ii v$ with $u,v \in \R$ and using
$u_{x_1} = v_{x_2}$, $u_{x_2} = -v_{x_1}$ gives, on $\{F \neq 0\}$,
\[
\abs{\nabla \abs{F}^2}^2
= (2 u u_{x_1} + 2 v v_{x_1})^2
+ (2 u u_{x_2} + 2 v v_{x_2})^2
= 4\,(u^2+v^2)(u_{x_1}^2 + v_{x_1}^2)
= 4\,\abs{F}^2 \abs{F'}^2.
\]
The zero set of a non-trivial entire function has Lebesgue measure zero,
so~\eqref{eq:bochner-pointwise} holds a.e.\ on~$\C$. Combining
Lemma~\ref{lem:weighted-equiv} (applied to $H = F^2 - F(0)^2 \in
\mathcal{F}^2_\ast(\C)$) with~\eqref{eq:bochner-pointwise},
\begin{equation}\label{eq:squaring-step1}
\norm{F^2 - F(0)^2}_{\mathcal{F}}^2
\;\asymp\;
\int_{\C} \abs{\nabla \abs{F}^2}^2\, \frac{\dd\gamma(z)}{1+\abs{z}^2}.
\end{equation}

\medskip
\noindent\emph{Step 2: First Bochner identity.}\;
Set $w(z) = (1+\abs{z}^2)^{-1}/\pi \cdot \ee^{-\abs{z}^2}$, so that the
right-hand side of \eqref{eq:squaring-step1} equals
$\int_{\C}\abs{\nabla\abs{F}^2}^2\, w\,\dd m$. Let $v$ be any real
harmonic function on $\C$ with $\abs{F}^2 - v \in L^2(\dd\gamma)$. Since
$\Delta v = 0$ and $\nabla \abs{F}^2 = \nabla(\abs{F}^2 - v) + \nabla v$,
two integrations by parts (boundary terms vanish thanks to the Gaussian
factor in $w$) yield
\begin{equation}\label{eq:bochner-1st}
\int_{\C} \abs{\nabla \abs{F}^2}^2\, w\, \dd m
\;=\;
-\int_{\C} (\abs{F}^2 - v)\, \Delta\abs{F}^2\, w\, \dd m
+ \tfrac{1}{2} \int_{\C} (\abs{F}^2 - v)^2\, \Delta w\, \dd m.
\end{equation}
A direct computation shows that there exists an absolute constant $C_0$
with
\begin{equation}\label{eq:weight-laplacian}
\abs{\Delta w(z)}
\;\le\;
\frac{C_0}{\pi}\, \frac{\ee^{-\abs{z}^2}}{1+\abs{z}^2}
\qquad\text{for all } z \in \C.
\end{equation}
Indeed, $w = (1/\pi)\ee^{-\abs{z}^2}/(1+\abs{z}^2)$ satisfies
$\abs{\nabla \log w} \le 2\abs{z} + 2\abs{z}/(1+\abs{z}^2)$ and a similar
bound on $\Delta\log w$, so combining gives
$\abs{\Delta w}/w \le C_0\,(1 + \abs{z}^2)$, which is
\eqref{eq:weight-laplacian} after multiplying through.
Substituting~\eqref{eq:weight-laplacian} into the second term
of~\eqref{eq:bochner-1st},
\begin{equation}\label{eq:bochner-1st-2nd-term}
\Bigl|\, \tfrac{1}{2}\int_{\C} (\abs{F}^2 - v)^2\, \Delta w\, \dd m \,\Bigr|
\;\le\;
\frac{C_0}{2}\, \bigl\|\, \abs{F}^2 - v \,\bigr\|_{L^2(\dd\gamma)}^2.
\end{equation}

\medskip
\noindent\emph{Step 3: Cauchy--Schwarz on the principal term.}\;
For the first term in~\eqref{eq:bochner-1st}, Cauchy--Schwarz applied
with the splitting $w = (\ee^{-\abs{z}^2/2}/\sqrt{\pi}) \cdot
(\ee^{-\abs{z}^2/2}/(\sqrt{\pi}(1+\abs{z}^2)))$ gives
\begin{equation}\label{eq:bochner-CS}
\Bigl|\,\int_{\C} (\abs{F}^2-v)\, \Delta\abs{F}^2\, w\, \dd m \,\Bigr|
\;\le\;
\bigl\|\,\abs{F}^2 - v\,\bigr\|_{L^2(\dd\gamma)}\, J^{1/2},
\end{equation}
where
\begin{equation}\label{eq:J-def}
J \;:=\; \int_{\C} \bigl(\Delta\abs{F}^2\bigr)^2\,
\frac{\dd\gamma(z)}{(1+\abs{z}^2)^2}.
\end{equation}

\medskip
\noindent\emph{Step 4: Bounding $J$ via a second Bochner identity.}\;
We now claim that
\begin{equation}\label{eq:J-bound}
J \;\le\; C_1\, \bigl\|\,\abs{F}^2 - v\,\bigr\|_{L^2(\dd\gamma)}^2
\end{equation}
for some absolute constant $C_1$. Since $v$ is real harmonic,
$\Delta(\abs{F}^2 - v) = \Delta\abs{F}^2$, so using the symmetry of the
Laplacian under integration by parts twice we may write
\[
J
= \int_{\C} \Delta(\abs{F}^2 - v)\, \Delta\abs{F}^2\,
\frac{\dd\gamma}{(1+\abs{z}^2)^2}
= \int_{\C} (\abs{F}^2 - v)\,
\Delta\!\left(\Delta\abs{F}^2\, w_2\right) \dd m,
\]
where $w_2(z) := (1/\pi)\ee^{-\abs{z}^2}/(1+\abs{z}^2)^2$. Expanding the
outer Laplacian by the product rule,
\[
\Delta\!\left(\Delta\abs{F}^2\, w_2\right)
= \Delta^2\abs{F}^2\, w_2
\;+\; 2\, \nabla(\Delta\abs{F}^2) \cdot \nabla w_2
\;+\; \Delta\abs{F}^2\, \Delta w_2,
\]
hence $J = I_1 + I_2 + I_3$ with the obvious notation. We bound each
$I_j$ in turn using:
\begin{itemize}[leftmargin=2em]
    \item the identity $\Delta\abs{F}^2 = 4\abs{F'}^2$ (which gives
    $\Delta^2\abs{F}^2 = 16\abs{F''}^2$ by another application of
    \eqref{eq:bochner-pointwise} to $F'$);
    \item the pointwise bound $\bigl|\nabla(\Delta\abs{F}^2)\bigr|
    = 4\,\abs{\nabla\abs{F'}^2} \le 8\,\abs{F'}\abs{F''}$
    (from~\eqref{eq:fnt-nabla-mod} applied to $F'$);
    \item the weight bounds $\abs{\nabla w_2} \le C_0'\,
    (1/\pi)\ee^{-\abs{z}^2}/(1+\abs{z}^2)^{3/2}$ and
    $\abs{\Delta w_2} \le C_0'\,
    (1/\pi)\ee^{-\abs{z}^2}/(1+\abs{z}^2)$, both of which are derived
    exactly as~\eqref{eq:weight-laplacian}.
\end{itemize}
Cauchy--Schwarz against $\abs{F}^2 - v$ then gives, with $\dd\gamma_k(z)
:= (1+\abs{z}^2)^{-k}\,\dd\gamma(z)$ for shorthand,
\begin{align*}
\abs{I_1}
&\;\le\; 16\,\bigl\|\,\abs{F}^2-v\,\bigr\|_{L^2(\dd\gamma)}\,
\Bigl(\int_{\C}\abs{F''}^4\,\dd\gamma_4\Bigr)^{\!1/2}, \\
\abs{I_2}
&\;\le\; C_2'\,\bigl\|\,\abs{F}^2-v\,\bigr\|_{L^2(\dd\gamma)}\,
\Bigl(\int_{\C}\abs{F'}^2\abs{F''}^2\,\dd\gamma_3\Bigr)^{\!1/2}, \\
\abs{I_3}
&\;\le\; C_2'\,\bigl\|\,\abs{F}^2-v\,\bigr\|_{L^2(\dd\gamma)}\,
\Bigl(\int_{\C}\abs{F'}^4\,\dd\gamma_2\Bigr)^{\!1/2}
= C_2'\,\bigl\|\,\abs{F}^2-v\,\bigr\|_{L^2(\dd\gamma)}\, J^{1/2}/4,
\end{align*}
the last identity using $\Delta\abs{F}^2 = 4\abs{F'}^2$ in the definition
\eqref{eq:J-def} of $J$.

\medskip

The two integrals appearing in $\abs{I_1}$ and $\abs{I_2}$ are controlled
by the integral defining $J$ via the following weighted Cauchy integral
estimate, whose proof we postpone to Subsection~\ref{sec:weighted-cauchy}.

\begin{lemma}\label{lem:weighted-cauchy}
There exists an absolute constant $C_3 > 0$ such that for every polynomial
$F \in \mathcal{F}^2(\C)$,
\begin{align}
\int_{\C} \abs{F''(z)}^4\, \dd\gamma_4(z)
&\;\le\; C_3 \int_{\C} \abs{F'(z)}^4\, \dd\gamma_2(z),
\label{eq:cauchy-l4}\\
\int_{\C} \abs{F'(z)}^2\abs{F''(z)}^2\, \dd\gamma_3(z)
&\;\le\; C_3 \int_{\C} \abs{F'(z)}^4\, \dd\gamma_2(z).
\label{eq:cauchy-mixed}
\end{align}
\end{lemma}

Granting Lemma~\ref{lem:weighted-cauchy} and recalling $\int\abs{F'}^4
\,\dd\gamma_2 = J/16$, we obtain
\[
\max\bigl(\abs{I_1},\abs{I_2},\abs{I_3}\bigr)
\;\le\;
C_4\,\bigl\|\,\abs{F}^2 - v\,\bigr\|_{L^2(\dd\gamma)}\, J^{1/2},
\]
hence $J \le 3 C_4\,\bigl\|\,\abs{F}^2-v\,\bigr\|_{L^2(\dd\gamma)}\,
J^{1/2}$, which gives~\eqref{eq:J-bound} with $C_1 = 9 C_4^2$.

\medskip
\noindent\emph{Step 5: Conclusion.}\;
Substituting \eqref{eq:bochner-1st-2nd-term},
\eqref{eq:bochner-CS} and \eqref{eq:J-bound}
into~\eqref{eq:bochner-1st} and combining with~\eqref{eq:squaring-step1},
\[
\norm{F^2 - F(0)^2}_{\mathcal{F}}^2
\;\le\;
C\,\bigl\|\,\abs{F}^2 - v\,\bigr\|_{L^2(\dd\gamma)}^2,
\]
which is Theorem~\ref{thm:squaring}. The minimum
in~\eqref{eq:squaring-min} over real constants is attained at
$c = \int\abs{F}^2\,\dd\gamma = \norm{F}_{\mathcal{F}}^2$, and any harmonic
function is admissible since the argument used $\Delta v = 0$ only.
\qed

\subsubsection{The weighted Cauchy integral lemma}\label{sec:weighted-cauchy}

The purpose of this subsubsection is to prove a quantitative form of
Lemma~\ref{lem:weighted-cauchy}. Throughout, we fix the cut-off radius
\begin{equation}\label{eq:R-choice}
R \;=\; 2,
\end{equation}
and split each integral into the bounded region $\{\abs{z}\le R\}$ and
its complement, using the Cauchy circle of radius~$1$ on the former
and the Cauchy circle of radius $1/\abs{z}$ on the latter.  We
shall track all constants explicitly along the way; the values
displayed below are chosen for legibility rather than sharpness, and
follow at once from the choice~\eqref{eq:R-choice}.

\medskip

We first state the quantified version we are about to prove.

\begin{lemma}[Quantitative weighted Cauchy integral lemma]\label{lem:weighted-cauchy-quant}
For every polynomial $F \in \mathcal{F}^2(\C)$,
\begin{align}
\int_{\C} \abs{F''(z)}^4\, \dd\gamma_4(z)
&\;\le\;
A_1\,\int_{\C} \abs{F'(z)}^4\, \dd\gamma_2(z),
\label{eq:cauchy-l4-quant}\\
\int_{\C} \abs{F'(z)}^2 \abs{F''(z)}^2\, \dd\gamma_3(z)
&\;\le\;
\sqrt{A_1}\;\int_{\C} \abs{F'(z)}^4\, \dd\gamma_2(z),
\label{eq:cauchy-mixed-quant}
\end{align}
with the explicit constant
\begin{equation}\label{eq:A1-explicit}
A_1
\;:=\;
\ee^{2R+1}\,\bigl(1+(R+1)^2\bigr)^{2}
\;+\;
\frac{\ee^{\,2+1/R^{2}}}{(1-1/R^{2})\,\bigl(1-2/(1+R^{2})\bigr)^{2}}
\;\le\; 1.5\times 10^{4}.
\end{equation}
\end{lemma}

\noindent
The constants $C_3$ and $C_3'$ in
Lemma~\ref{lem:weighted-cauchy} may therefore be taken to be
$A_1$ and $\sqrt{A_1}$ respectively.

\medskip

The proof of Lemma~\ref{lem:weighted-cauchy-quant} occupies the rest of
this subsubsection.  We prepare three independent ingredients:
the pointwise Cauchy bound for $F''$ in terms of a circular average of
$F'$ (Lemma~\ref{lem:pointwise-cauchy} below); a quantitative comparison
of weights under the substitution $w = z + r_z\ee^{\ii\theta}$
(Lemma~\ref{lem:weight-shift}); and the deduction
of~\eqref{eq:cauchy-mixed-quant} from~\eqref{eq:cauchy-l4-quant} via
Cauchy--Schwarz.

\begin{lemma}[Pointwise Cauchy bound]\label{lem:pointwise-cauchy}
Let $F \colon \C \to \C$ be entire and let $r > 0$.  Then for every
$z \in \C$,
\begin{equation}\label{eq:pointwise-cauchy}
\abs{F''(z)}^{4}
\;\le\;
\frac{1}{2\pi\, r^{4}} \int_{0}^{2\pi}
\abs{F'(z + r\,\ee^{\ii\theta})}^{4}\, \dd\theta.
\end{equation}
\end{lemma}

\begin{proof}
The Cauchy integral formula on the circle $\{\abs{w-z} = r\}$ gives
\[
F''(z)
\;=\;
\frac{1}{2\pi\ii}\int_{\abs{w-z}=r} \frac{F'(w)}{(w-z)^{2}}\,\dd w
\;=\;
\frac{1}{2\pi r}\int_{0}^{2\pi} F'(z + r\,\ee^{\ii\theta})\,
\ee^{-\ii\theta}\,\dd\theta.
\]
Taking moduli and applying Jensen's inequality with the convex function
$t \mapsto t^{4}$,
\[
\abs{F''(z)}^{4}
\;\le\;
\Bigl(\frac{1}{2\pi r}\int_{0}^{2\pi}
\abs{F'(z + r\,\ee^{\ii\theta})}\,\dd\theta\Bigr)^{4}
\;\le\;
\frac{1}{(2\pi r)^{4}}\,(2\pi)^{3}
\int_{0}^{2\pi}\abs{F'(z + r\,\ee^{\ii\theta})}^{4}\,\dd\theta,
\]
which simplifies to~\eqref{eq:pointwise-cauchy}.
\end{proof}

\begin{lemma}[Quantitative weight shift]\label{lem:weight-shift}
Let $R \ge 2$.
\begin{enumerate}[\rm(i)]
\item For every $z \in \C$ with $\abs{z}\le R$, every $\theta \in
[0,2\pi)$, and $w := z + \ee^{\ii\theta}$,
\begin{equation}\label{eq:weight-shift-bd}
\frac{\ee^{-\abs{z}^{2}}}{(1+\abs{z}^{2})^{4}}
\;\le\;
\ee^{2R+1}\,\bigl(1+(R+1)^{2}\bigr)^{2}\;
\frac{\ee^{-\abs{w}^{2}}}{(1+\abs{w}^{2})^{2}}.
\end{equation}
\item For every $z \in \C$ with $\abs{z} > R$, every $\theta \in
[0,2\pi)$, and $w := z + \ee^{\ii\theta}/\abs{z}$, the map $z \mapsto
w(z,\theta)$ is a $C^{\infty}$-diffeomorphism of $\{\abs{z}>R\}$ onto
its image, with absolute Jacobian determinant satisfying
\[
\bigl\lvert J(z)\bigr\rvert \;\ge\; 1 - \frac{1}{R^{2}},
\]
and the weight comparison
\begin{equation}\label{eq:weight-shift-out}
\frac{\abs{z}^{4}\,\ee^{-\abs{z}^{2}}}{(1+\abs{z}^{2})^{4}}
\;\le\;
\frac{\ee^{\,2 + 1/R^{2}}}{\bigl(1-2/(1+R^{2})\bigr)^{2}}\;
\frac{\ee^{-\abs{w}^{2}}}{(1+\abs{w}^{2})^{2}}.
\end{equation}
\end{enumerate}
\end{lemma}

\begin{proof}
\emph{(i)\,Bounded region.}\;
Since $w = z+\ee^{\ii\theta}$ we have $\abs{w}\le \abs{z}+1\le R+1$ and
\[
\abs{z}^{2} - \abs{w}^{2}
\;=\; -\,2\,\operatorname{Re}\!\bigl(z\,\ee^{-\ii\theta}\bigr) - 1
\;\le\; 2\abs{z} - 1
\;\le\; 2R - 1
\;<\; 2R + 1.
\]
Hence
$\ee^{-\abs{z}^{2}} \le \ee^{2R+1}\,\ee^{-\abs{w}^{2}}$. Moreover
$(1+\abs{z}^{2})^{-4} \le 1$, while $(1+\abs{w}^{2})^{2} \le
(1+(R+1)^{2})^{2}$. Multiplying these three inequalities
yields~\eqref{eq:weight-shift-bd}.

\medskip

\emph{(ii)\,Outer region: Jacobian.}\;
Writing $z = x+\ii y$ and $w = u + \ii v$ with $u = x +
\cos\theta/\abs{z}$ and $v = y + \sin\theta/\abs{z}$, a direct
computation using $\partial_{x}\abs{z} = x/\abs{z}$ and
$\partial_{y}\abs{z} = y/\abs{z}$ gives
\[
\det\!\begin{pmatrix} \partial_{x} u & \partial_{y} u \\
\partial_{x} v & \partial_{y} v \end{pmatrix}
\;=\;
1 \;-\; \frac{\operatorname{Re}\bigl(\ee^{-\ii\theta}\,z\bigr)}{\abs{z}^{3}}.
\]
Since $\abs{\operatorname{Re}(\ee^{-\ii\theta}z)} \le \abs{z}$, the
absolute Jacobian determinant is bounded below by
$1 - 1/\abs{z}^{2} \ge 1 - 1/R^{2}$, which is strictly positive for
$R\ge 2$.  Hence the map is a $C^{\infty}$-local diffeomorphism, and
since $\abs{w-z} = 1/\abs{z}$ tends to $0$ as $\abs{z}\to\infty$, it is
also injective on $\{\abs{z}>R\}$ provided $R \ge 2$ (which we have
fixed).

\medskip

\emph{(ii)\,Outer region: weights.}\;
We have $\abs{w}^{2} = \abs{z}^{2} +
2\,\operatorname{Re}(z\,\ee^{-\ii\theta})/\abs{z} + 1/\abs{z}^{2}$, so
\[
\abs{z}^{2} - \abs{w}^{2}
\;=\; -\,2\,\frac{\operatorname{Re}\bigl(z\,\ee^{-\ii\theta}\bigr)}{\abs{z}}
- \frac{1}{\abs{z}^{2}}
\;\in\; [-2 - 1/R^{2},\; 2 + 1/R^{2}],
\]
giving $\ee^{\abs{z}^{2}-\abs{w}^{2}} \le \ee^{\,2 + 1/R^{2}}$.
For the polynomial weight, $\abs{w}\ge \abs{z} - 1/\abs{z}$, hence on
$\{\abs{z} > R\}$,
\[
\frac{1+\abs{w}^{2}}{1+\abs{z}^{2}}
\;\ge\;
\frac{1+(\abs{z}-1/\abs{z})^{2}}{1+\abs{z}^{2}}
\;=\;
1 \;-\; \frac{2 - 1/\abs{z}^{2}}{1+\abs{z}^{2}}
\;\ge\;
1 - \frac{2}{1+R^{2}}.
\]
Combining and using $\abs{z}^{4}/(1+\abs{z}^{2})^{2} \le 1$,
\[
\frac{\abs{z}^{4}\,\ee^{-\abs{z}^{2}}}{(1+\abs{z}^{2})^{4}}
\;=\;
\frac{\abs{z}^{4}}{(1+\abs{z}^{2})^{2}}\cdot
\frac{\ee^{-\abs{z}^{2}}}{(1+\abs{z}^{2})^{2}}
\;\le\;
\frac{\ee^{\,\abs{z}^{2}-\abs{w}^{2}}\,(1+\abs{w}^{2})^{2}}
{(1+\abs{z}^{2})^{2}}\cdot
\frac{\ee^{-\abs{w}^{2}}}{(1+\abs{w}^{2})^{2}}
\;\le\;
\frac{\ee^{\,2 + 1/R^{2}}}{\bigl(1-2/(1+R^{2})\bigr)^{2}}\cdot
\frac{\ee^{-\abs{w}^{2}}}{(1+\abs{w}^{2})^{2}},
\]
which is~\eqref{eq:weight-shift-out}.
\end{proof}

\begin{proof}[Proof of Lemma~\ref{lem:weighted-cauchy-quant}]
We first establish~\eqref{eq:cauchy-l4-quant}, then deduce
\eqref{eq:cauchy-mixed-quant} via Cauchy--Schwarz between the two
$L^{4}$ estimates.

\medskip
\noindent\emph{Step 1: Splitting and pointwise bound.}\;
Let $r_{z} := 1$ if $\abs{z}\le R$ and $r_{z} := 1/\abs{z}$
if $\abs{z}>R$, with $R = 2$ as fixed
in~\eqref{eq:R-choice}. Lemma~\ref{lem:pointwise-cauchy} applied
with $r = r_{z}$ gives, for every $z \in \C$,
\[
\abs{F''(z)}^{4}\,\dd\gamma_{4}(z)
\;\le\;
\frac{1}{2\pi\, r_{z}^{4}}\int_{0}^{2\pi}
\abs{F'(z + r_{z}\,\ee^{\ii\theta})}^{4}\, \dd\theta\,
\frac{\ee^{-\abs{z}^{2}}}{\pi\,(1+\abs{z}^{2})^{4}}\,\dd m(z).
\]
We integrate this inequality first over $\{\abs{z}\le R\}$, then over
$\{\abs{z}>R\}$, exchanging the orders of $z$- and
$\theta$-integration by Fubini's theorem.

\medskip
\noindent\emph{Step 2: Inner contribution} ($\abs{z}\le R$).
On this region $r_{z} = 1$, so
\[
\frac{1}{r_{z}^{4}}\,\frac{\ee^{-\abs{z}^{2}}}{(1+\abs{z}^{2})^{4}}
\;\le\;
\ee^{2R+1}\,(1+(R+1)^{2})^{2}\,\frac{\ee^{-\abs{w}^{2}}}{(1+\abs{w}^{2})^{2}}
\]
by Lemma~\ref{lem:weight-shift}\,(i), with $w = z + \ee^{\ii\theta}$.
The map $z \mapsto z + \ee^{\ii\theta}$ is a translation, hence has unit
Jacobian, and sends $\{\abs{z}\le R\}$ into $\{\abs{w}\le R+1\}$.
Therefore
\[
\int_{\abs{z}\le R}\!\frac{\dd m(z)}{2\pi^{2} r_{z}^{4}}\!
\int_{0}^{2\pi}\!\!\abs{F'(z + \ee^{\ii\theta})}^{4}\,
\frac{\ee^{-\abs{z}^{2}}}{(1+\abs{z}^{2})^{4}}\,\dd\theta
\;\le\;
\ee^{2R+1}\,(1+(R+1)^{2})^{2}
\int_{\C}\!\abs{F'(w)}^{4}\,\dd\gamma_{2}(w),
\]
where the right-hand side has been enlarged from $\{\abs{w}\le R+1\}$
to all of $\C$ since the integrand is nonnegative.

\medskip
\noindent\emph{Step 3: Outer contribution} ($\abs{z}>R$).
On this region $r_{z} = 1/\abs{z}$, so
\[
\frac{1}{r_{z}^{4}}\,\frac{\ee^{-\abs{z}^{2}}}{(1+\abs{z}^{2})^{4}}
\;=\;
\frac{\abs{z}^{4}\,\ee^{-\abs{z}^{2}}}{(1+\abs{z}^{2})^{4}}.
\]
Lemma~\ref{lem:weight-shift}\,(ii) bounds this expression, with
$w = z + \ee^{\ii\theta}/\abs{z}$, by
\[
\frac{\ee^{\,2 + 1/R^{2}}}{(1-2/(1+R^{2}))^{2}}\,
\frac{\ee^{-\abs{w}^{2}}}{(1+\abs{w}^{2})^{2}},
\]
and the map $z\mapsto w$ is a diffeomorphism with absolute Jacobian
determinant $\ge 1 - 1/R^{2}$, so the change of variables produces an
extra factor at most $(1-1/R^{2})^{-1}$. Hence
\[
\int_{\abs{z}>R}\frac{\dd m(z)}{2\pi^{2}r_{z}^{4}}
\int_{0}^{2\pi}\!\!\abs{F'(z+\ee^{\ii\theta}/\abs{z})}^{4}
\frac{\ee^{-\abs{z}^{2}}}{(1+\abs{z}^{2})^{4}}\,\dd\theta
\;\le\;
\frac{\ee^{\,2+1/R^{2}}}
{(1-1/R^{2})\,\bigl(1-2/(1+R^{2})\bigr)^{2}}
\int_{\C}\abs{F'(w)}^{4}\,\dd\gamma_{2}(w).
\]
Note that the factor $\abs{z}^{4}$ produced by $r_{z}^{-4}$ has already
been absorbed into~\eqref{eq:weight-shift-out} via the elementary
bound $\abs{z}^{4}/(1+\abs{z}^{2})^{2}\le 1$.

\medskip
\noindent\emph{Step 4: Adding the two contributions.}\;
Combining Steps~2 and~3,
\[
\int_{\C}\abs{F''(z)}^{4}\,\dd\gamma_{4}(z)
\;\le\;
A_{1}\,\int_{\C}\abs{F'(w)}^{4}\,\dd\gamma_{2}(w),
\]
with $A_{1}$ given by~\eqref{eq:A1-explicit}. For $R = 2$,
\[
\ee^{2R+1}\,(1+(R+1)^{2})^{2}
\;=\; \ee^{5}\cdot 100 \;\le\; 14842,
\qquad
\frac{\ee^{2+1/R^{2}}}{(1-1/R^{2})\,(1-2/(1+R^{2}))^{2}}
\;=\;
\frac{\ee^{9/4}}{(3/4)\cdot(3/5)^{2}}
\;\le\; 36,
\]
giving $A_{1}\le 14842 + 36 < 1.5\times 10^{4}$, which
is~\eqref{eq:cauchy-l4-quant}.

\medskip
\noindent\emph{Step 5: Mixed estimate.}\;
For~\eqref{eq:cauchy-mixed-quant}, write the integrand of
$\int_{\C}\abs{F'}^{2}\abs{F''}^{2}\,\dd\gamma_{3}$ as
\[
\bigl(\abs{F'(z)}^{2}\,(1+\abs{z}^{2})^{-1}\bigr)\,
\bigl(\abs{F''(z)}^{2}\,(1+\abs{z}^{2})^{-2}\bigr)\,
\frac{\ee^{-\abs{z}^{2}}}{\pi}\,\dd m(z),
\]
and apply Cauchy--Schwarz against the Gaussian density:
\[
\int_{\C}\abs{F'}^{2}\abs{F''}^{2}\,\dd\gamma_{3}
\;\le\;
\Bigl(\int_{\C}\abs{F'}^{4}\,\dd\gamma_{2}\Bigr)^{1/2}\,
\Bigl(\int_{\C}\abs{F''}^{4}\,\dd\gamma_{4}\Bigr)^{1/2}.
\]
Substituting~\eqref{eq:cauchy-l4-quant} into the second factor,
\[
\int_{\C}\abs{F'}^{2}\abs{F''}^{2}\,\dd\gamma_{3}
\;\le\;
\Bigl(\int_{\C}\abs{F'}^{4}\,\dd\gamma_{2}\Bigr)^{1/2}
\sqrt{A_{1}}\,\Bigl(\int_{\C}\abs{F'}^{4}\,\dd\gamma_{2}\Bigr)^{1/2}
\;=\;
\sqrt{A_{1}}\int_{\C}\abs{F'}^{4}\,\dd\gamma_{2},
\]
which is~\eqref{eq:cauchy-mixed-quant}. With $A_{1}\le 1.5\times
10^{4}$, the constant in~\eqref{eq:cauchy-mixed-quant} is
$\sqrt{A_{1}}\le 123$.
\end{proof}

\begin{remark}\label{rem:cauchy-sharper}
The constants in~\eqref{eq:A1-explicit} are far from sharp. Two
optimizations are immediate. First, the inner-region constant
$\ee^{2R+1}(1+(R+1)^{2})^{2}$ is the only place where $R$ enters
exponentially; it can be reduced by replacing the unit-radius Cauchy
circle on $\{\abs{z}\le R\}$ with a circle of radius $\rho < 1$, at the
cost of a factor $\rho^{-4}$ in front, after which one optimises in
$\rho$.  Second, the polynomial bound
$\abs{z}^{4}/(1+\abs{z}^{2})^{2}\le 1$ used in Step~3 can be sharpened
to $1 - 2/(1+\abs{z}^{2})^{2}$, which integrates against the Gaussian
to a term that is again exponentially small in $R$. Carrying out both
optimisations and choosing $R$ near $1$ leads to constants of order
$10^{2}$, but we have not pursued these refinements as they play no
quantitative role in the application to Theorem~\ref{thm:squaring}.
\end{remark}

\subsubsection{An improved Cauchy--Schwarz inequality}\label{sec:improved-CS}

The bound \eqref{eq:squaring-min} translates into a quantitative
strict inequality in Cauchy--Schwarz between $\operatorname{Re}(G)$ and
$\abs{G}^2$ in $L^2(\dd\gamma)$, valid for entire functions vanishing at
the origin.

\begin{theorem}\label{thm:improved-cs}
There exists an absolute constant $c_\sharp > 0$ such that for every
non-zero polynomial $G \in \mathcal{F}^2_\ast(\C)$,
\begin{equation}\label{eq:improved-cs}
\Bigl\|\, \frac{\abs{G}^2}{\norm{G}_{L^4(\dd\gamma)}^2}
\;-\; \frac{\operatorname{Re}(G)}{\norm{\operatorname{Re}(G)}_{L^2(\dd\gamma)}}
\Bigr\|_{L^2(\dd\gamma)}
\;\ge\; c_\sharp.
\end{equation}
Equivalently, there exists $\eta_\sharp \in (0,1)$ such that for every
$G \in \mathcal{F}^2_\ast(\C)$,
\begin{equation}\label{eq:cs-strict}
\Bigl| \int_{\C}\operatorname{Re}(G)\,\abs{G}^2\,\dd\gamma\Bigr|
\;\le\;
\eta_\sharp\, \norm{\operatorname{Re}(G)}_{L^2(\dd\gamma)}\,
\norm{\,\abs{G}^2\,}_{L^2(\dd\gamma)}.
\end{equation}
\end{theorem}

\begin{proof}
We may assume $G$ is normalized so that $\norm{G}_{L^4(\dd\gamma)} = 1$,
i.e.\ $\norm{G^2}_{\mathcal{F}} = 1$ (since $G^2$ is entire and the Fock
norm coincides with the $L^2(\dd\gamma)$ norm on $\mathcal{F}^2(\C)$).
The hypothesis $G(0)=0$ gives $G^2(0)=0$, so $G^2 \in
\mathcal{F}^2_\ast(\C)$ and the orthogonality of monomials yields
\begin{equation}\label{eq:F2-norm-1}
\min_{\alpha \in \C} \norm{G^2 - \alpha}_{L^2(\dd\gamma)}
= \norm{G^2}_{L^2(\dd\gamma)}
= 1.
\end{equation}
On the other hand, $\operatorname{Re}(G) =
\tfrac{1}{2}(G + \overline{G})$ is harmonic, so the function
$v := \alpha\, \operatorname{Re}(G)$ is harmonic for every $\alpha \in \R$.
Choosing
$\alpha = \norm{\operatorname{Re}(G)}_{L^2(\dd\gamma)}^{-1}$, the bound
\eqref{eq:squaring-min} of Theorem~\ref{thm:squaring} (with
$F = G$ and $F(0) = 0$) gives
\[
1
\;=\; \norm{G^2}_{\mathcal{F}}
\;\le\; C\,\Bigl\|\, \abs{G}^2
\;-\; \frac{\operatorname{Re}(G)}{\norm{\operatorname{Re}(G)}_{L^2(\dd\gamma)}}\,
\Bigr\|_{L^2(\dd\gamma)}.
\]
Hence \eqref{eq:improved-cs} holds with $c_\sharp = 1/C$.

\medskip

For \eqref{eq:cs-strict}, let
$\rho := \langle \abs{G}^2, \operatorname{Re}(G)\rangle_{L^2(\dd\gamma)} /
\bigl(\norm{\,\abs{G}^2\,}_{L^2(\dd\gamma)}\,
\norm{\operatorname{Re}(G)}_{L^2(\dd\gamma)}\bigr)$ be the normalized
inner product. Squaring~\eqref{eq:improved-cs} yields
$2(1 - \rho) \ge c_\sharp^2$, so $\rho \le 1 - c_\sharp^2/2$. Replacing
$G$ by $-G$ flips the sign of $\operatorname{Re}(G)$ but not $\abs{G}^2$,
so the same bound applies to $-\rho$. We conclude
$\abs{\rho} \le 1 - c_\sharp^2/2 =: \eta_\sharp < 1$, which is
\eqref{eq:cs-strict}.
\end{proof}

\begin{remark}\label{rem:cs-implies-orth}
Theorem~\ref{thm:improved-cs} provides a route to the orthogonal
coercivity bound of Theorem~\ref{thm:orthog} in the small-perturbation
regime. Indeed, expanding pointwise
$\abs{1+U}^2 - 1 = 2\operatorname{Re}(U) + \abs{U}^2$
and integrating against $\dd\gamma$,
\[
\int_{\C}\bigl(\abs{1+U}^2-1\bigr)^2\,\dd\gamma
= 4\norm{\operatorname{Re}(U)}_{L^2(\dd\gamma)}^2
+ 4\!\int_{\C}\!\operatorname{Re}(U)\,\abs{U}^2\,\dd\gamma
+ \norm{\,\abs{U}^2\,}_{L^2(\dd\gamma)}^2.
\]
Inserting~\eqref{eq:cs-strict} and completing the square gives
\[
\int_{\C}\bigl(\abs{1+U}^2-1\bigr)^2\,\dd\gamma
\;\ge\;
(1-\eta_\sharp)\Bigl(4\norm{\operatorname{Re}(U)}_{L^2(\dd\gamma)}^2
+ \norm{\,\abs{U}^2\,}_{L^2(\dd\gamma)}^2\Bigr).
\]
Since $U \in \mathcal{F}^2_\ast(\C)$ has only positive frequencies, the
real and imaginary parts of $U$ have equal $L^2(\dd\gamma)$ norms, hence
$\norm{U}_{\mathcal{F}}^2 = 2\norm{\operatorname{Re}(U)}_{L^2(\dd\gamma)}^2$,
and we obtain
\[
\int_{\C}\bigl(\abs{1+U}^2-1\bigr)^2\,\dd\gamma
\;\ge\;
2(1-\eta_\sharp)\,\norm{U}_{\mathcal{F}}^2.
\]
Factoring $\abs{1+U}^2 - 1 = (\abs{1+U}+1)\,(\abs{1+U}-1)$ and using
$\abs{1+U}+1 \le 2 + \abs{U}$, the left-hand side is bounded by a
universal constant times $\norm{\rho \circ U}_{\mathcal{F}}^2$ provided
$\norm{U}_{\mathcal{F}}$ is sufficiently small; this recovers
Theorem~\ref{thm:orthog} in the local-local regime of
Theorem~\ref{thm:local}.
\end{remark}

\subsubsection{A convexity bound for moment integrals}\label{sec:convexity}

We close this subsection with a separate but related observation,
recording a sharp convexity estimate for the family of moment integrals
naturally associated with the deficit $f := \abs{1+rG}^2 - 1$.

\begin{proposition}\label{prop:convexity}
Let $\mu$ be any probability measure on $\C$ and let $f : \C \to \R$ be a
$\mu$-measurable function with $f \ge -1$ a.e. Define, for $s \in (0,1)$,
\[
g(s) \;:=\; \int_{\C} \frac{f(z)^2}{\bigl(1 + s\,f(z)\bigr)^{3/2}}\,
\dd\mu(z).
\]
Suppose $g$ is finite on $[0,1)$, i.e.\ $f^2 \in L^1(\mu)$ (so that
$g(0) = \int f^2\,\dd\mu < \infty$). Then $g \in C^2((0,1))$ with
\begin{equation}\label{eq:diff-ineq}
\tfrac{5}{3}\,(g'(s))^2 \;\le\; g(s)\, g''(s)
\qquad\text{for all } s \in (0,1).
\end{equation}
In particular, if $g'(0) < 0$, then for every $t \in (0,1)$,
\begin{equation}\label{eq:g-lower}
g(t)
\;\ge\;
g(0)\, \Bigl(1 + \tfrac{2}{3}\,t\,\frac{\abs{g'(0)}}{g(0)}\Bigr)^{-3/2}.
\end{equation}
\end{proposition}

\begin{proof}
Differentiating under the integral sign,
\[
g'(s) = -\tfrac{3}{2}\!\int_{\C}\frac{f^3\,\dd\mu}{(1+sf)^{5/2}},
\qquad
g''(s) = \tfrac{15}{4}\!\int_{\C}\frac{f^4\,\dd\mu}{(1+sf)^{7/2}}.
\]
The bound $f \ge -1$ ensures $1 + sf \ge 1-s > 0$ for $s \in (0,1)$, and
the integrals are finite by H\"older against $f^2/(1+sf)^{3/2} \in
L^1(\mu)$. The Cauchy--Schwarz inequality applied to the decomposition
\[
\frac{f^3}{(1+sf)^{5/2}}
= \frac{f}{(1+sf)^{3/4}} \cdot \frac{f^2}{(1+sf)^{7/4}}
\]
gives, after squaring and integrating,
\[
\bigl(g'(s)\bigr)^2
= \tfrac{9}{4}\Bigl(\!\int_{\C}\!\frac{f^3\,\dd\mu}{(1+sf)^{5/2}}\Bigr)^{\!2}
\;\le\; \tfrac{9}{4}\!\int_{\C}\!\frac{f^2\,\dd\mu}{(1+sf)^{3/2}}
\!\int_{\C}\!\frac{f^4\,\dd\mu}{(1+sf)^{7/2}}
= \tfrac{9}{4}\cdot g(s)\cdot \tfrac{4}{15}\,g''(s),
\]
which is~\eqref{eq:diff-ineq}.

\medskip

Assume now $g'(0) < 0$, and let $I = (0,\alpha_0)$ be the maximal interval
on which $g'(s) < 0$. Set $h(s) := g'(s)/g(s)$. Rearranging
\eqref{eq:diff-ineq} as $\tfrac{5}{3}\,(g'/g)^2 \le g''/g$ and using
$h'(s) = g''(s)/g(s) - (g'(s)/g(s))^2$, we obtain
\[
\tfrac{2}{3}\,h(s)^2
\;\le\; h'(s)
\quad\Longleftrightarrow\quad
-\tfrac{2}{3}
\;\ge\; -\frac{h'(s)}{h(s)^2}
= \partial_s\!\Bigl(\frac{1}{h(s)}\Bigr).
\]
Integrating over $[0,t] \subset I$ and using $h(0) = g'(0)/g(0) < 0$,
\[
\frac{1}{h(t)} - \frac{1}{h(0)}
\;\le\; -\tfrac{2}{3}\,t,
\qquad
\text{i.e.}\qquad
\frac{g(t)}{g'(t)}
\;\le\; \frac{g(0)}{g'(0)} - \tfrac{2}{3}\,t.
\]
Since $g'(t),g'(0) < 0$ on $I$, multiplying by $g'(t)\,g'(0)>0$ and
rearranging yields
\[
\frac{g'(t)}{g(t)}
\;\ge\; \frac{g'(0)}{g(0) - \tfrac{2}{3}\,t\,g'(0)}.
\]
Finally, integrating once more from $0$ to $t$ via the substitution
$v = -\tfrac{2}{3}g'(0)\,s$,
\[
\log\!\Bigl(\frac{g(t)}{g(0)}\Bigr)
\;\ge\;
\int_0^t \frac{g'(0)\,\dd s}{g(0) - \tfrac{2}{3}\,g'(0)\,s}
= -\tfrac{3}{2}\,
\log\!\Bigl(1 - \tfrac{2}{3}\,t\,\frac{g'(0)}{g(0)}\Bigr),
\]
which is~\eqref{eq:g-lower}. The case $g'(0) \ge 0$ is even easier:
$g$ is non-decreasing on $[0, \alpha_0)$ by~\eqref{eq:diff-ineq}, so
$g(t) \ge g(0)$ on this interval.
\end{proof}

Specializing $\mu = \gamma$ and $f = \abs{1+rG}^2 - 1$ for $G \in
\mathcal{F}^2(\C)$ and $r > 0$, the bound \eqref{eq:g-lower} gives a
quantitative coercive control of the family
$s \mapsto \int_{\C} f^2\,(1+sf)^{-3/2}\,\dd\gamma$
in terms of its initial value $\int f^2\,\dd\gamma =
\norm{\,\abs{1+rG}^2 - 1\,}_{L^2(\dd\gamma)}^2$, complementing the
$L^2$-stability provided by Theorem~\ref{thm:squaring}.
